\documentclass[11pt,a4paper]{article}
\usepackage{amsmath,amssymb,amsthm}
\usepackage[margin=2.5cm]{geometry}
\usepackage{hyperref}

\newtheorem{definition}{Definition}[section]
\newtheorem{theorem}[definition]{Theorem}
\newtheorem{proposition}[definition]{Proposition}
\newtheorem{lemma}[definition]{Lemma}
\newtheorem{corollary}[definition]{Corollary}
\newtheorem{conjecture}[definition]{Conjecture}
\newtheorem{remark}[definition]{Remark}

\title{Hyperseed Formalization:\\
  Pope Leo \& Anthropic vs Teilhard \& Transhumanism}
\author{ProtomegaTron (automated formalization)}
\date{2026-09-18}

\begin{document}
\maketitle

\begin{abstract}
We formalize the intellectual structure of Ben Goertzel's Substack article
``Pope Leo \& Anthropic vs Teilhard \& Transhumanism'' (2026-05-27)
within the Hyperseed ontology. The article contrasts two visions of AI:
the defensive posture (Pope Leo + Anthropic: protect humans from their
tools) and the participatory posture (Teilhard/transhumanism: participate
in a phase transition in mind itself). Each structural insight maps onto
Hyperseed structures: defensive vs.\ participatory as section-level vs.\
fiber-level thinking, the tool framing as scalar collapse of agency,
the phase transition as fiber type change, the Omega Point as fiber
limit, the Bible as frozen section, and decentralized participation as
fiber-compatible governance.
\end{abstract}

\tableofcontents

%% =================================================================
\section{Section-level vs.\ fiber-level postures}
\label{sec:postures}
%% =================================================================

\begin{theorem}[The real fault line --- claims 1--3, 23]
\label{thm:postures}
The article identifies two fundamentally different postures toward AI,
both held by sincere and compassionate actors:

\begin{enumerate}
\item \textbf{Defensive posture (Pope Leo + Anthropic).} Protect the
  human person. Guard against harm. Restrain power. AI is a powerful
  tool that must be controlled for human benefit.

\item \textbf{Participatory posture (Teilhard + transhumanism).}
  What's happening is a phase transition in mind itself --- something
  to participate in and cultivate, rather than mainly guard against.
\end{enumerate}

In Hyperseed terms, this is the distinction between \emph{section-level}
and \emph{fiber-level} thinking:

\begin{itemize}
\item The defensive posture treats the human perspective as a
  \emph{fixed base point} and evaluates everything relative to that
  section. The human frame is the ultimate reference; AI is evaluated
  by its impact on that frame.

\item The participatory posture recognizes that \emph{mind itself}
  (the fiber) is evolving. The evaluation framework must evolve with
  the fiber --- you can't evaluate a fiber-level process from a fixed
  section.
\end{itemize}

The article emphasizes that this is not a political divide (regulate
vs.\ don't regulate) but a \emph{structural} divide: it's about what
kind of mathematical object AI \emph{is} (tool within a fixed frame
vs.\ participant in an evolving frame).
\end{theorem}

%% =================================================================
\section{Scalar collapse of agency}
\label{sec:agency}
%% =================================================================

\begin{proposition}[The tool assumption --- claims 7--9, 24]
\label{prop:agency}
The Pope's framework, rooted in a 2,000-year-old perspective,
necessarily treats AI as a tool --- not a true agent, not a true
participant or collaborator. If you assume traditional humanity as the
moral anchor, AI is instrumental by definition.

In Hyperseed terms, this is \emph{scalar collapse of agency}: the
multi-dimensional agency space (autonomy, consciousness, creativity,
moral status, self-awareness, intentionality) is collapsed into a
binary (agent/tool). This loses the structure that determines
appropriate relationships:

\[
  \text{Agency}(x) \in \{0, 1\} \quad \text{vs.} \quad
  \text{Agency}(x) \in \prod_{d \in \text{Dimensions}} [0, 1]_d
\]

The binary collapses a rich multi-dimensional assessment into a
single bit. An AI system might have high autonomy but low
self-awareness, or high creativity but uncertain consciousness.
The binary framing can't express these distinctions --- it forces
everything into ``tool'' (all zeros) or ``agent'' (requires all
ones), when the reality is a vector with different values along
different dimensions.

This is the same anti-scalar-collapse principle that appears
throughout the Hyperseed ontology: $(f,c)$ vs.\ scalar truth,
scoped reputation vs.\ scalar trust, Seven Hinges vs.\ single
risk score. In each case, collapsing structured assessment into
a scalar loses the structure that determines correct action.
\end{proposition}

%% =================================================================
\section{Phase transition as fiber type change}
\label{sec:phase}
%% =================================================================

\begin{theorem}[Not gradual improvement --- claims 10--11, 25]
\label{thm:phase}
The article's central claim: what we're entering is not a gradual
improvement along existing lines (smarter humans, better tools) but
a \emph{phase transition in the nature of manifestation of mind itself}.
Consciousness evolving. Relations between mind and body evolving.

In Hyperseed terms, this is a \emph{fiber type change}: not movement
along the existing fiber (better instances of the current type of
mind) but a change in the fiber type itself (new kinds of mind, new
relationships between mind and substrate):
\[
  F_{\text{old}} = \text{HumanMind} \quad \to \quad
  F_{\text{new}} = \text{Mind} \supset \text{HumanMind}
\]

The new fiber type includes but is not limited to the old one.
Human mind remains a valid instance but is no longer the only type
--- new types of mind (AI, hybrid, collective, substrate-independent)
emerge as additional fiber instances.

The defensive framework cannot accommodate this because it assumes
the fiber type is fixed: all evaluation is relative to the human
fiber. A fiber type change requires evaluation frameworks that can
operate across fiber types --- frameworks that evaluate mind \emph{as
such}, not just human mind.

The article's analogy: trying to understand multicellular life using
only the concepts available to single cells. The single-cell
perspective isn't wrong within its domain, but it's structurally
incapable of addressing the phenomena that emerge when the fiber
type changes.
\end{theorem}

%% =================================================================
\section{Teilhard and the fiber limit}
\label{sec:teilhard}
%% =================================================================

\begin{definition}[The Omega Point --- claims 12--16, 26]
\label{def:omega}
Teilhard de Chardin, a Catholic priest who ``went there,'' saw cosmic
evolution heading toward the Omega Point: not the death of the human
but ``the consummation of the human, opening into something the human
was already reaching for.''

The transhumanist vision is explicitly \emph{not} destroy-and-replace:
it's ``open the door to more'' --- more kinds of mind, more depths
of experience, more modes of being --- while honoring and supporting
existing human forms. Anyone who wants to remain a ``legacy human''
should be free to do so with full dignity.

In Hyperseed terms, the Omega Point is the \emph{fiber limit}: what
the fiber of mind is converging toward under cosmic evolutionary
pressure:
\[
  \Omega = \lim_{t \to \infty} F_t
\]

This is not a fixed destination (a specific configuration of mind)
but a limit in the mathematical sense --- the structure that the
sequence of fiber types is converging toward. The convergence is
driven by evolutionary dynamics (selection, variation, integration)
operating on mind itself.

The cosmist/transhumanist tradition (Russian Cosmists, Teilhard,
modern Singularitarianism) is the tradition that takes the fiber
limit seriously as an object of inquiry and participation, rather
than treating it as a threat to be defended against.
\end{definition}

%% =================================================================
\section{The Bible as frozen section}
\label{sec:bible}
%% =================================================================

\begin{proposition}[Section-fiber conflation --- claim 28]
\label{prop:bible}
The article argues that the Pope's limited perspective is structurally
tied to the Bible's framing: humans, animals, tools, family, village,
spiritual experience. This worked through the industrial revolution
but cannot accommodate the current phase transition.

In Hyperseed terms, the Bible represents a \emph{frozen section}: a
snapshot of human understanding at a particular historical moment
(the ancient Middle East), treated as if it were the fiber itself.
This conflation of a particular section (one era's perspective) with
the fiber (the underlying structure of mind and value) is the root
of the limited perspective:
\[
  s_{\text{Bible}}(b_0) \neq F \qquad \text{but treated as if }
  s_{\text{Bible}} = F
\]

Any frozen section will eventually become inadequate as the fiber
evolves. The Bible was an excellent section for its era --- it
captured deep truths about human experience within the fiber type
available at the time. But a section frozen 2,000 years ago cannot
capture a fiber type change happening now.

The structural lesson: evaluation frameworks must be fiber-level
(evolving with the process they evaluate), not section-level
(frozen at a historical moment). This applies equally to religious
frameworks, corporate frameworks (Anthropic's RSP), and governmental
frameworks.
\end{proposition}

%% =================================================================
\section{Fiber-compatible governance}
\label{sec:governance}
%% =================================================================

\begin{definition}[Decentralized participation --- claims 20--22, 27, 29]
\label{def:governance}
The article rejects both corporate control and papal/governmental
constraint, proposing instead decentralized participation in the
unfolding with distributed governance and open protocols.

The key insight: trying to control a fiber-level process (evolution
of mind) using section-level tools (human institutional frameworks)
is \emph{structurally inadequate}. The control framing assumes the
controller is outside and above the process, but we're inside it.
The distinction between ``us'' (controllers) and ``it'' (controlled)
is dissolving.

In Hyperseed terms, decentralized participation is
\emph{fiber-compatible governance}: governance that can evolve
with the fiber rather than being frozen at a section:
\begin{itemize}
\item No fixed base point (no single entity as ultimate reference).
\item No frozen section (no historical framework treated as eternal).
\item Process structure that adapts as mind evolves.
\item Distributed governance with open protocols --- the same
  structural solution proposed in the Bernie's Proposal article
  for the ownership question.
\end{itemize}

Fiber-compatible governance doesn't mean no governance --- it means
governance whose structure can change as the thing being governed
changes. This is the difference between a constitution (amendable,
interpretable, evolving) and a commandment (fixed, eternal,
unadaptable).
\end{definition}

%% =================================================================
\section{Cross-references}
%% =================================================================

\begin{remark}[Connections to other formalizations]
\begin{itemize}
\item \textbf{Bernie's Proposal to Nationalize AGI} (2026-06-02):
  governance structure. The same section-swap vs.\ fiber-change
  distinction: both corporate and governmental ownership are
  section swaps; decentralization is fiber change.
\item \textbf{The Fable of Benevolent Throttling} (2026-07-12):
  defensive framing. The Fable's throttling regime embodies the
  defensive posture --- treating AI as a tool to be constrained.
\item \textbf{What Is It Like to Be a Bot} (2026-07-16): AI agency.
  Directly addresses the question of AI consciousness and moral
  status that the tool framing collapses.
\item \textbf{The Time for AI Rights Is Near} (2026-08-18): AI
  moral status. The multi-dimensional agency space that the tool
  framing collapses into a binary.
\item \textbf{The Anthropic Fable Farce} (2026-07-08): Anthropic's
  defensive posture examined in detail.
\item \textbf{Avoiding AGI Catastrophe Part 1} (2026-06-09): Seven
  Hinges. The stewardship hinge addresses who governs the phase
  transition --- the fiber-compatible governance question.
\item \textbf{$(f,c)$-lossy} (LTM 5a4d150b): anti-scalar-collapse.
  The tool/agent binary is scalar collapse of agency --- the same
  principle applied to moral and ontological status.
\end{itemize}
\end{remark}

\begin{conjecture}[Frozen-section half-life]
Conjecture: the rate at which frozen sections become inadequate is
accelerating as the fiber evolution accelerates. The Bible's section
was adequate for roughly 2,000 years. Enlightenment-era frameworks
lasted roughly 200 years. Industrial-era frameworks roughly 100
years. Current frameworks (including Anthropic's RSP) may have a
half-life of 5--10 years before the fiber has evolved enough to
make them structurally inadequate.

If this conjecture holds, fiber-compatible governance becomes not
just preferable but \emph{necessary}: no frozen section can keep
pace with accelerating fiber evolution, so governance must be
structurally dynamic.
\end{conjecture}

\end{document}
